\documentclass[a4paper,11pt]{article}
\usepackage[utf8]{inputenc}
\usepackage[T1]{fontenc}
\usepackage[margin=2.5cm]{geometry}
\usepackage{hyperref}
\usepackage{booktabs}
\usepackage{longtable}
\usepackage{array}
\usepackage{enumitem}
\usepackage{xcolor}
\usepackage{fancyhdr}
\usepackage{titlesec}

% ── Colours ──────────────────────────────────────────────────────────────────
\definecolor{accentblue}{RGB}{13,110,253}
\definecolor{lightgray}{RGB}{245,245,245}
\definecolor{warnred}{RGB}{185,28,28}

% ── Header / footer ──────────────────────────────────────────────────────────
\pagestyle{fancy}
\fancyhf{}
\rhead{\small Class Fund Manager}
\lhead{\small Teacher \& Treasurer Manual}
\rfoot{\small Page \thepage}
\lfoot{\small \today}

% ── Section formatting ────────────────────────────────────────────────────────
\titleformat{\section}{\large\bfseries\color{accentblue}}{}{0em}{}[\titlerule]
\titleformat{\subsection}{\normalsize\bfseries}{}{0em}{}

% ── Hyperlink colours ─────────────────────────────────────────────────────────
\hypersetup{
    colorlinks=true,
    linkcolor=accentblue,
    urlcolor=accentblue
}

% ─────────────────────────────────────────────────────────────────────────────
\title{\textbf{Class Fund Manager}\\[0.4em]
       \large Teacher \& Treasurer Manual}
\author{}
\date{\today}
% ─────────────────────────────────────────────────────────────────────────────

\begin{document}

\maketitle
\thispagestyle{empty}
\vspace{1em}

\noindent
This manual covers every task a \textbf{teacher}, \textbf{Full Treasurer}, or
\textbf{Bookkeeper} will need to perform in the Class Fund Manager.
The application runs at:
\begin{center}
    \url{https://two025-wt-prj-dembinny.onrender.com}
\end{center}

\tableofcontents
\newpage

% ─────────────────────────────────────────────────────────────────────────────
\section{Roles and Permissions}
% ─────────────────────────────────────────────────────────────────────────────

The system has six user roles.
The table below shows which actions each role may perform.

\begin{longtable}{>{\bfseries}p{4.5cm} p{8.5cm}}
\toprule
Role & Capabilities \\
\midrule
\endhead
System Admin    & Full access to every class and every feature.
                  Creates and manages class accounts. \\
Treasurer Full  & Full treasurer capabilities for their assigned class:
                  payment requests, transactions, expenses, student management,
                  bank account visibility. \\
Treasurer Accountant & Can manage payment requests and view the bank account,
                  but cannot log expenses or manage students. \\
Treasurer Bookkeeper & Can log expenses and record cash transactions,
                  but cannot create payment requests or manage students. \\
Student         & Views their own assigned payment requests, pays via QR,
                  sees their own transaction history. \\
Parent          & Views the same payment data as their linked child. \\
\bottomrule
\end{longtable}

\noindent
\textbf{Fund Groups} are custom permission sets you create inside the Admin
Panel (see Section~\ref{sec:admin}). A Fund Group acts like a named role for
your class — useful when you want a subset of treasurer abilities for class
representatives without giving full treasurer access.

% ─────────────────────────────────────────────────────────────────────────────
\section{Logging In}
% ─────────────────────────────────────────────────────────────────────────────

Navigate to \url{https://two025-wt-prj-dembinny.onrender.com/login/} and enter
your \textbf{username} and \textbf{password}.
After a successful login you are redirected to the Dashboard appropriate for
your role.

If you forget your password, contact your System Admin — there is no
self-service password reset in the current version.

% ─────────────────────────────────────────────────────────────────────────────
\section{Dashboard Overview}
% ─────────────────────────────────────────────────────────────────────────────

The Treasurer Dashboard is your home screen. It shows:

\begin{itemize}[noitemsep]
    \item \textbf{Fund Balance} — total collected minus total expenses
          (Collected: sum of all confirmed transactions;
           Spent: sum of all logged expenses).
    \item \textbf{Payment Requests} — a list of active requests with how many
          students have paid (\emph{N/Total}).
    \item \textbf{Recent Expenses} — the most recent expense records.
    \item \textbf{Quick links} — navigate to Expenses, Import, or Admin Panel.
\end{itemize}

\noindent
\textcolor{warnred}{\textbf{Important:}} The ``Fund Balance'' displayed in the
top navigation bar respects the \texttt{hide\_fund\_balance} flag on each
student's account. If a student has this flag set, they cannot see the fund
total.

% ─────────────────────────────────────────────────────────────────────────────
\section{Bulk Importing Students (CSV)}
\label{sec:import}
% ─────────────────────────────────────────────────────────────────────────────

The fastest way to populate your class roster is to upload a CSV file.
Only users with \textbf{Treasurer Full} or \textbf{System Admin} privileges
can access the importer.

\subsection{How to navigate to the importer}

\begin{enumerate}[noitemsep]
    \item From the Dashboard click \textbf{Import Students}.
    \item Select your \textbf{Class} from the dropdown.
    \item Upload the \texttt{.csv} file and click \textbf{Preview Import}.
    \item Review the preview table — green rows will be created,
          red rows have errors and will be skipped.
    \item Click \textbf{Confirm Import} to apply all valid rows.
\end{enumerate}

\subsection{CSV format and column reference}

The file must be \textbf{UTF-8 encoded} and include a \textbf{header row}.
Column names are case-insensitive and leading/trailing spaces are ignored.

\bigskip
\begin{longtable}{>{\ttfamily}p{4cm} p{2cm} p{7.5cm}}
\toprule
\normalfont Column & Required? & Notes \\
\midrule
\endhead
username           & \textit{see note} & Login name used for authentication. \\
email              & \textit{see note} & Student email. When no username is given,
                                         the local part (before @) becomes the username. \\
first\_name        & \textit{see note} & First name. Used together with \texttt{last\_name}
                                         to derive a username when neither username nor email
                                         is provided (e.g.\ ``Jan Novák'' → ``jnovak''). \\
last\_name         & Optional  & Last name. Stored on the user account. \\
role               & Optional  & One of: \texttt{student}, \texttt{parent},
                                  \texttt{treasurer\_full}, \texttt{treasurer\_accountant},
                                  \texttt{treasurer\_bookkeeper},
                                  or any \textbf{Fund Group name} you have created for this class.
                                  Defaults to \texttt{student} when omitted. \\
variable\_symbol   & Optional  & Payment identifier (5 digits).
                                  Auto-generated from the class prefix if blank. \\
parent\_email      & Optional  & If provided, a Parent account is created (or found by email)
                                  and linked to this student. \\
parent\_first\_name & Optional & First name stored on the Parent account. \\
parent\_last\_name  & Optional & Last name stored on the Parent account. \\
password           & Optional  & Sets the initial password. A random 12-character
                                  password is generated if blank.
                                  Communicate the password to the user separately. \\
\bottomrule
\end{longtable}

\noindent
\textbf{Identity rule:} Every row must supply \emph{at least one} of:
\texttt{username}, \texttt{email}, or both \texttt{first\_name} and
\texttt{last\_name}. Rows that have none of these are rejected.

\subsection{Minimal CSV example}

\begin{verbatim}
first_name,last_name,parent_email
Jan,Novák,jan.novak.otec@example.cz
Marie,Horáková,horak.maminka@example.cz
\end{verbatim}

\subsection{Full CSV example with all columns}

\begin{verbatim}
username,email,first_name,last_name,role,variable_symbol,
    parent_email,parent_first_name,parent_last_name,password
jnovak,jan.novak@skola.cz,Jan,Novák,student,04001,
    jan.novak.otec@example.cz,Petr,Novák,tajne123
mhora,,Marie,Horáková,student,,
    horak.maminka@example.cz,Eva,Horáková,
\end{verbatim}

\noindent
\textcolor{warnred}{\textbf{Note:}} Lines are wrapped here for readability.
In the actual file each row must be a single unbroken line.

\subsection{What happens on import?}

For each valid row the system:
\begin{enumerate}[noitemsep]
    \item Looks up an existing user by \texttt{username} → \texttt{email}
          → first+last name. Creates a new user if none is found.
    \item Sets the user's \texttt{role}. If the role matches a \textbf{Fund Group}
          name, the user is created as a \texttt{TREASURER\_BOOKKEEPER} and added
          to that group.
    \item Generates a \textbf{Variable Symbol} if not supplied
          (format: \texttt{<2-digit class prefix><3-digit sequence>}).
    \item Creates or links the \textbf{Parent account} if \texttt{parent\_email}
          is present.
    \item Records the import in the \texttt{ImportJob} log for audit purposes.
\end{enumerate}

% ─────────────────────────────────────────────────────────────────────────────
\section{Creating a Payment Request}
% ─────────────────────────────────────────────────────────────────────────────

Payment Requests tell students and parents what they owe and by when.

\begin{enumerate}[noitemsep]
    \item Navigate to \textbf{Finances → Payment Requests → New Request}.
    \item Fill in:
        \begin{itemize}[noitemsep]
            \item \textbf{Title} — e.g.\ "School Trip to Prague".
            \item \textbf{Amount} — decimal, e.g.\ \texttt{450.00} CZK.
            \item \textbf{Due Date} — students see this deadline.
            \item \textbf{Assign to all} — tick to assign every active class
                  student automatically.
                  Untick to choose specific students from the list.
            \item \textbf{Variable Symbol / Specific Symbol} — optional;
                  helps match bank transfers if used.
        \end{itemize}
    \item Click \textbf{Save}. The request is immediately visible on every
          assigned student's dashboard.
\end{enumerate}

\noindent
A SPAYD QR code is automatically generated for each student using their
personal Variable Symbol. See Section~\ref{sec:qr}.

% ─────────────────────────────────────────────────────────────────────────────
\section{Recording Transactions}
% ─────────────────────────────────────────────────────────────────────────────

\subsection{Manual cash payment}

Use this when a student brings cash in person.

\begin{enumerate}[noitemsep]
    \item Open \textbf{Finances → Payment Requests} and click the relevant
          request.
    \item Find the student in the list and click \textbf{Record Cash Payment}.
    \item Enter the amount received (may be partial).
    \item A transaction is created with status \textbf{PENDING}.
    \item Click \textbf{Confirm} to mark it as received.
          The student's dashboard immediately shows it as paid.
\end{enumerate}

\subsection{Bank transfer (planned)}

The system is designed to import bank transactions automatically from
a Fio Bank account. This feature is \textbf{not yet live} — matching is
planned to use the student's Variable Symbol.
Until then, all bank transfers must be confirmed manually as cash payments.

\subsection{Transaction statuses}

\begin{center}
\begin{tabular}{ll}
\toprule
Status & Meaning \\
\midrule
PENDING   & Recorded but not yet confirmed by a treasurer. \\
CONFIRMED & Payment accepted; contributes to fund balance. \\
REJECTED  & Declined (wrong amount, bounced, etc.). \\
\bottomrule
\end{tabular}
\end{center}

% ─────────────────────────────────────────────────────────────────────────────
\section{Logging an Expense}
% ─────────────────────────────────────────────────────────────────────────────

Expenses reduce the displayed fund balance and give parents transparency
into where the class money is spent.

\begin{enumerate}[noitemsep]
    \item Navigate to \textbf{Finances → Expenses → New Expense}.
    \item Fill in:
        \begin{itemize}[noitemsep]
            \item \textbf{Title} — what was purchased.
            \item \textbf{Amount} — decimal, e.g.\ \texttt{2500.00} CZK.
            \item \textbf{Category} — one of:
                  Trip, Supplies, Food, Decoration, Donation, Other.
            \item \textbf{Spent at} — date of the purchase.
            \item \textbf{Publish} — tick to make the expense visible
                  to students and parents in their dashboard.
        \end{itemize}
    \item Click \textbf{Save}.
\end{enumerate}

\noindent
Expenses are shown in the treasurer dashboard; published expenses also
appear in the student/parent view.

% ─────────────────────────────────────────────────────────────────────────────
\section{SPAYD QR Codes}
\label{sec:qr}
% ─────────────────────────────────────────────────────────────────────────────

Each payment request page displays a QR code in the
\textbf{SPAYD format} — the Czech banking standard for mobile payments.

When a parent scans the QR code with their banking app:
\begin{itemize}[noitemsep]
    \item The class bank account number is pre-filled.
    \item The payment amount is pre-filled.
    \item The student's \textbf{Variable Symbol} is pre-filled.
\end{itemize}

\noindent
This means the bank transfer can be matched to the correct student
automatically once bank import is activated.
Parents never need to type account numbers or symbols manually.

\noindent
The QR codes are generated server-side and embedded directly in the page
as Base64 PNG images — no external service is contacted.

% ─────────────────────────────────────────────────────────────────────────────
\section{Managing Classes and Users (Admin Panel)}
\label{sec:admin}
% ─────────────────────────────────────────────────────────────────────────────

The Django Admin Panel at \url{.../admin/} is accessible to \textbf{System Admins} only.

\subsection{Creating a new class}

\begin{enumerate}[noitemsep]
    \item \textbf{Admin → Accounts → School Classes → Add School Class}.
    \item Set \textbf{Name} (e.g.\ "4A"), \textbf{Teacher}, \textbf{School Year},
          and \textbf{VS Prefix} (a 2-digit number uniquely identifying the class
          for Variable Symbol generation — e.g.\ \texttt{04}).
    \item Save.
\end{enumerate}

\subsection{Creating a Fund Group}

Fund Groups let you assign custom treasurer sub-roles to users
(e.g.\ a class representative who can log expenses but not manage payment
requests).

\begin{enumerate}[noitemsep]
    \item \textbf{Admin → Accounts → Fund Groups → Add Fund Group}.
    \item Select the \textbf{Class}.
    \item Enter the \textbf{Name} — this name can be used as the
          \texttt{role} column value in CSV imports.
    \item Tick the \textbf{permission flags} that apply to this group:
        \begin{itemize}[noitemsep]
            \item \textit{Can log expenses}
            \item \textit{Can manage payment requests}
            \item \textit{Can view bank account}
            \item \textit{Can manage students}
        \end{itemize}
    \item Save.
\end{enumerate}

\subsection{Adding or editing a Class Membership}

A \textbf{Class Membership} links a user to a class and optionally
assigns them a Fund Group.

\begin{enumerate}[noitemsep]
    \item \textbf{Admin → Accounts → Class Memberships → Add}.
    \item Choose \textbf{User}, \textbf{School Class}, and optionally a
          \textbf{Fund Group}.
    \item Save.
\end{enumerate}

\subsection{Managing the Bank Account record}

\begin{enumerate}[noitemsep]
    \item \textbf{Admin → Finances → Bank Accounts → Add}.
    \item Link to a \textbf{School Class} (one active bank account per class).
    \item Enter \textbf{Account Number}, \textbf{IBAN}, \textbf{BIC},
          \textbf{Bank Name}, and \textbf{Owner Name}.
    \item Tick \textbf{Is Active} and Save.
\end{enumerate}

% ─────────────────────────────────────────────────────────────────────────────
\section{Viewing the Notification Log}
% ─────────────────────────────────────────────────────────────────────────────

Every email sent by the system (welcome messages, payment reminders,
receipts) is stored in the \textbf{Notification Log}:

\begin{enumerate}[noitemsep]
    \item \textbf{Admin → Communications → Notification Logs}.
    \item Use the filters on the right to narrow by \textit{type} or
          \textit{recipient}.
    \item Each record shows the recipient, type, timestamp, and the
          rendered message body.
\end{enumerate}

\noindent
This log is read-only — it is an audit trail, not a re-send queue.

% ─────────────────────────────────────────────────────────────────────────────
\section{Common Troubleshooting}
% ─────────────────────────────────────────────────────────────────────────────

\begin{longtable}{p{5.5cm} p{8cm}}
\toprule
Problem & Solution \\
\midrule
\endhead
A student cannot log in.
    & Check their username in Admin → Accounts → Users.
      Reset the password from the user's edit page. \\
Import shows all rows as errors.
    & Ensure the file is saved as \textbf{UTF-8} (not UTF-16 or Windows-1250).
      Each row needs at least username, email, or first+last name. \\
Fund balance looks wrong.
    & Confirm that all cash transactions have been set to
      \textbf{CONFIRMED} status.
      PENDING transactions are not counted in the balance. \\
QR code does not appear.
    & The class must have an active Bank Account record (see Section~\ref{sec:admin}).
      Without an IBAN, no QR can be generated. \\
A student has no Variable Symbol.
    & Open \textbf{Admin → Accounts → Student Profiles}, find the student,
      and click \textbf{Regenerate VS}. Make sure the class has a
      non-empty \texttt{vs\_prefix}. \\
The site shows ``502 Bad Gateway'' on Render.
    & The free Render plan spins down after inactivity.
      Wait 30–60 seconds for the dyno to wake up, then refresh. \\
\bottomrule
\end{longtable}

% ─────────────────────────────────────────────────────────────────────────────
\section{Contact and Support}
% ─────────────────────────────────────────────────────────────────────────────

This project is maintained as a school web-technologies assignment.
For bug reports or feature requests open an issue on GitHub:

\begin{center}
    \url{https://github.com/gyarab/2025_wt_prj_dembinny}
\end{center}

\end{document}
